\thispagestyle{empty}

% --- PHẦN TIÊU ĐỀ (Mô phỏng theo mẫu ảnh bạn gửi) ---
\noindent
\vspace{0.2cm}
\begin{center}
    \textbf{\LARGE TÓM TẮT}
\end{center}
\vspace{1cm} % Khoảng cách xuống nội dung

% --- NỘI DUNG TÓM TẮT ---

Đồ án chuyên ngành này tập trung vào việc nghiên cứu, thiết kế và xây dựng \textbf{Hệ thống Quản lý Thư viện Trực tuyến tích hợp AI}, nhằm giải quyết các hạn chế của các hệ thống thư viện truyền thống như khả năng tìm kiếm kém hiệu quả và thiếu sự cá nhân hóa. Hệ thống được phát triển với mục tiêu số hóa toàn diện quy trình vận hành thư viện, đồng thời ứng dụng trí tuệ nhân tạo để nâng cao trải nghiệm đọc và khám phá tri thức của người dùng.

Về mặt kỹ thuật, hệ thống được xây dựng dựa trên kiến trúc \textbf{Modular Monolith} với Clean Architecture + Domain-Driven Design cho từng module, đảm bảo tính bảo mật và khả năng mở rộng. Backend sử dụng \textbf{Spring Boot (Java)} để xử lý các nghiệp vụ phức tạp, kết hợp với Frontend \textbf{React + Vite} mang lại giao diện hiện đại, tốc độ build nhanh và triển khai đơn giản qua NGINX. Các chức năng chính bao gồm: quản lý tài nguyên sách, quản lý độc giả, quy trình mượn trả tự động, kiểm soát phí phạt và hệ thống báo cáo thống kê trực quan cho thủ thư.

Điểm đột phá của đồ án là việc tích hợp sâu các công nghệ AI/NLP tiên tiến. Cụ thể, hệ thống ứng dụng mô hình \textbf{Sentence-BERT (SBERT)} để thực hiện tìm kiếm ngữ nghĩa (Semantic Search), giúp xử lý các truy vấn tiếng Việt phức tạp mà không phụ thuộc vào từ khóa chính xác. Song song đó, hệ thống gợi ý sách sử dụng phương pháp lai (\textbf{Hybrid Approach}), kết hợp giữa Lọc cộng tác (Collaborative Filtering) và Lọc dựa trên nội dung (Content-Based), giúp đưa ra các đề xuất sách phù hợp nhất dựa trên lịch sử và hành vi của từng độc giả.

Kết quả của đồ án là một hệ thống phần mềm hoàn chỉnh, có giao diện thân thiện, đáp ứng đầy đủ các yêu cầu nghiệp vụ của một thư viện đại học. Hệ thống đã chứng minh được tính khả thi trong việc kết hợp giữa quản lý truyền thống và công nghệ thông minh, mang lại giá trị thực tiễn cao trong việc thúc đẩy văn hóa đọc và chuyển đổi số trong môi trường giáo dục.
